\begin{exercise}[The Newtonian tidal tensor]{MIT 8.962, Problem Set 7}
For a weak static field with \(g_{00}=-(1+2\Phi)\), compute
\(R_{0i0j}\) to leading order.  Use geodesic deviation to recover the
Newtonian relative-acceleration equation and interpret the vacuum trace.
\end{exercise}

\begin{exercise}[Volume of a small geodesic ball]{Cambridge Part III, General Relativity examples}
In Riemann normal coordinates prove
\[
\sqrt{|g|}=1-\frac16R_{\mu\nu}(p)x^\mu x^\nu+\order(x^3).
\]
For a small \(n\)-dimensional Riemannian geodesic ball, determine the first
curvature correction to its Euclidean volume.
\end{exercise}

\begin{exercise}[Gauss curvature from embedding data]{Tong, Cambridge Part III Sheet 2}
Derive the Gauss equation for a hypersurface in a flat ambient space.  Apply
it to a round two-sphere of radius \(a\) and recover its intrinsic Riemann
tensor and scalar curvature without computing Christoffel symbols.
\end{exercise}

\begin{exercise}[Codazzi and the momentum constraint]{MIT 8.962, Problem Set 10}
For a spacelike hypersurface with unit normal \(n^\mu\) and extrinsic
curvature \(K_{ij}\), derive the contracted Codazzi equation
\[
D_j(K^{ij}-h^{ij}K)=G_{\mu\nu}n^\mu h^{\nu i}.
\]
Explain its role when the Einstein equation holds.
\end{exercise}

\begin{exercise}[Jacobi fields and conjugate points]{Cambridge Part III, General Relativity examples}
Find the transverse Jacobi fields along a great circle on a two-sphere of
radius \(a\).  Locate the first conjugate point and explain why a great-circle
segment ceases to minimize distance beyond it.
\end{exercise}

\begin{exercise}[Conformal transformation of scalar curvature]{Cambridge Part III, 2017 exam}
In \(n\) dimensions, derive the transformation of \(R\) under
\(\widetilde g_{\mu\nu}=e^{2\omega}g_{\mu\nu}\).  Specialize to \(n=2\) and
relate the result to Gaussian curvature.
\end{exercise}

\begin{exercise}[A coordinate singularity versus a curvature singularity]{McGreevy, Problem Set 7}
For Schwarzschild spacetime, evaluate the Ricci scalar, Ricci-square, and
Kretschmann scalar.  Use them to distinguish the surfaces \(r=2M\) and
\(r=0\), and state what scalar invariants cannot by themselves establish.
\end{exercise}

\begin{exercise}[Electric and magnetic parts of Weyl curvature]{Cambridge Part III, General Relativity examples}
Relative to a unit timelike vector \(u^\mu\), define the electric and magnetic
parts of the Weyl tensor.  Show that each is spatial, symmetric, and trace
free, count their components, and reconstruct the ten independent components
of the Weyl tensor.
\end{exercise}
